\section{Five canonical regimes of the conservative dynamics}\label{sec:cases}

\begin{table*}
\caption{Five conservative regimes. Parameters in SI units ($k$ in N\,m$^{-1}$, $L$ in m, $M$ in kg, $g$ in m\,s$^{-2}$); every case starts at $\theta_0=0$ with angular velocity $\dot\theta_0$ (rad\,s$^{-1}$). The unit of time is $t_c=(L/g)^{1/2}$; $E$ is the mechanical energy, $E_{\mathrm{sep}}$ the separatrix energy, Eqs.~\eqref{eq:Esep_neg} and \eqref{eq:Esep}, and $T$ the exact period, Eq.~\eqref{eq:period}, of the libration or of one revolution.}
\label{tab:cases}
\begin{ruledtabular}
\begin{tabular}{cccccccccccl}
Case & $k$ & $L$ & $M$ & $g$ & $\dot\theta_0$ & $\mu$ & $t_c$ (s) & $E/MgL$ & $E_{\mathrm{sep}}/MgL$ & $T$ (s) & Motion\\
\hline
1 & 2   & 2 & 2 & 2   & 0.1 & $0$     & 1.000 & 1.005 & 1.000 & 16.389 & rotation\\
2 & 2   & 3 & 4 & 1   & 0.1 & $+0.50$ & 1.732 & 1.015 & 1.083 & 16.032 & libration, amplitude 0.252 rad\\
3 & 1   & 4 & 5 & 10  & 0.1 & $-0.92$ & 0.632 & 1.002 & 1.000 & 6.272  & rotation\\
4 & 0.1 & 4 & 2 & 0.1 & 0.4 & $+1.00$ & 6.325 & 4.200 & 1.250 & 14.838 & rotation\\
5 & 0.1 & 3 & 5 & 0.1 & 0.4 & $-0.40$ & 5.477 & 3.400 & 1.000 & 13.742 & rotation\\
\end{tabular}
\end{ruledtabular}
\end{table*}

\begin{figure*}
\includegraphics[width=\PIARlayout{0.92\textwidth}{\textwidth}]{figures/fig_cases.pdf}
\caption{Five conservative regimes of Table~\ref{tab:cases}, one case per row. Columns: (a) potential energy along the trajectory as a function of the angle (colored), potential landscape (gray), and total energy $E$ (dashed); (b) kinetic energy as a function of the angular velocity (colored) on the parabola $K=\frac12ML^{2}\dot\theta^{2}$ (gray); (c) phase diagram with the separatrix, Eq.~\eqref{eq:sep_curve} (dashed), centers (dots), and saddles (crosses); (d) position, i.e., the angle, as a function of time over 100~s. For rotations the angle is unwrapped. All curves are drawn from the dense output of DOP853 at 25\,001 instants. Full-size versions of the five rows are given in Figs.~S2--S6 of the supplementary material.}
\label{fig:cases}
\end{figure*}

\subsection{Parameter sets and numerical method}\label{sec:cases_method}

The predictions of Sec.~\ref{sec:model} are illustrated with the five parameter sets listed in Table~\ref{tab:cases}. They cover the three stability classes of the upright position, $kL=Mg$ (case 1), $kL>Mg$ (cases 2 and 4), and $kL<Mg$ (cases 3 and 5), combined with two initial speeds, and their time scales $t_c$ differ by an order of magnitude, from $0.63$ to $6.3$~s, which illustrates the similarity property of Sec.~\ref{sec:units}. Every case starts at the upright position, $\theta_0=0$, with a small angular velocity $\dot\theta_0$ of $0.1$ or $0.4\ \mathrm{rad\,s^{-1}}$. This protocol probes the property that matters most in applications, namely the response of an upright rod to a small kick. The parameters are expressed in SI units, and $g$ is treated as an effective gravitational acceleration: $g=10\ \mathrm{m\,s^{-2}}$ is standard gravity rounded, and the smaller values correspond to a pendulum that moves on a plane inclined by an angle $\alpha$ to the horizontal, for which $g_{\mathrm{eff}}=g\sin\alpha$. Because the dynamics in natural units depends only on $\mu$ (Sec.~\ref{sec:units}), the conclusions apply to every parameter set with the same $\mu$ and the same dimensionless kick $\dot\theta_0t_c$, that is, the same $E/MgL$.

The equation of motion was integrated with the explicit Runge--Kutta method DOP853, of order 8 with error estimators of orders 5 and 3~\cite{HairerNW1993} and built on the coefficients of Prince and Dormand~\cite{PrinceDormand1981}, as implemented in SciPy~\cite{Virtanen2020}, with relative and absolute tolerances of $10^{-12}$ and $10^{-13}$. Each trajectory was then evaluated at 25\,001 equally spaced instants over $100$~s through the continuous extension of order seven that accompanies the method, and all curves of Figs.~\ref{fig:cases} and \ref{fig:portraits} are drawn from this dense output. The mechanical energy, Eq.~\eqref{eq:H}, is conserved to within $8.1\times10^{-11}\,MgL$ in all cases, and the periods measured on the simulated trajectories agree with the quadrature of Eq.~\eqref{eq:period} to a relative accuracy better than $10^{-10}$. As an independent check, the five cases were integrated with the velocity Verlet scheme~\cite{Verlet1967,Swope1982} at a step of $T/2000$ with compensated summation~\cite{Kahan1965}. Because the scheme is symplectic, it conserves a modified Hamiltonian $\tilde H=H+O(h^{2})$, with $h$ the step, up to exponentially small errors over exponentially long times, so that its energy error stays bounded instead of drifting~\cite{BenettinGiorgilli1994,HairerLW2006}; compensated summation keeps the round-off error, which accumulates as a random walk (Brouwer's law~\cite{Brouwer1937}), far below this level. Its energy error remains below $5\times10^{-5}\,MgL$, and its angles agree with the DOP853 trajectories to within $3\times10^{-4}$~rad over the $100$~s.

\subsection{Energy landscapes, phase diagrams, and trajectories}\label{sec:cases_results}

Figure~\ref{fig:cases} presents each case through four complementary views. The first column shows the potential energy along the trajectory, $U(\theta(t))$ as a function of $\theta$, together with the total energy $E$; the part of the landscape that the rod explores, and the margin between $E$ and the peaks of $U$, decide the type of motion. The second column shows the kinetic energy, Eq.~\eqref{eq:K}, as a function of the angular velocity. The curve is always a segment of the parabola $K=\frac12ML^{2}\dot\theta^{2}$, and its extent reveals whether the velocity vanishes at turning points, as in a libration, or never does, as in a rotation. The third column is the phase diagram $(\theta,\dot\theta)$ with the separatrix of Eq.~\eqref{eq:sep_curve}, and the fourth column is the angle as a function of time. For rotations the angle is shown unwrapped, so that one revolution corresponds to an advance of $2\pi$. The cases are discussed from bound to fast motion: the libration (case 2), the two rotations that graze the separatrix (cases 1 and 3), and the two fast rotations (cases 4 and 5).

The upright position is stable in cases 2 and 4, but only in case 2 ($\mu=0.5$) does the rod remain in the upright well. The tilted saddles lie at $\theta^{*}=48.2^{\circ}$, and the energy supplied by the kick, $E=1.015\,MgL$, stays below the top of the barrier, $U(\theta^{*})=1.083\,MgL$. The rod therefore librates inside the upright well [Fig.~\ref{fig:cases}(a2)], with an amplitude of $0.252$~rad ($14.4^{\circ}$), along a closed orbit nested inside the lens-shaped separatrix that joins the two saddles [Fig.~\ref{fig:cases}(c2)]. The motion is nearly harmonic [Fig.~\ref{fig:cases}(d2)], but its period, $16.032$~s, exceeds the linear value $2\pi/\omega_0=15.391$~s by $4.2\%$, i.e., its frequency is lower by $4.0\%$. This lengthening is the signature of the softening cubic term of the normal form, Eq.~\eqref{eq:normal_form}. To first order in the squared amplitude $a^{2}$, the Lindstedt--Poincar\'e method~\cite{Strogatz2015} gives the frequency
\begin{equation}
\omega \simeq \omega_0\left[1-\frac{(3+4\mu)\,a^{2}}{16\,\mu}\right],
\label{eq:lindstedt}
\end{equation}
which predicts, for $a=0.2516$~rad, a frequency shift of $3.96\%$ against the exact $4.00\%$, that is, a period lengthening of $4.1\%$ against the exact $4.2\%$ (Appendix~\ref{app:derivations}).

Case 1 sits exactly at the critical point, $kL=Mg$. The upright position has no linear stiffness, which is the ideal of an isolator, but the quartic maximum $U\simeq MgL(1-\theta^{4}/8)$ makes it unstable, and the kick, which raises the energy only $0.5\%$ above the separatrix, suffices to make the rod rotate. The degenerate maximum produces a striking critical slowing down. Because the potential is flat near the top, the rod creeps across the upright position and then swings rapidly through the hanging one, so that it spends $71\%$ of the time within $45^{\circ}$ of the upright position. The angle advances in a staircase [Fig.~\ref{fig:cases}(d1)], and the orbit hugs the separatrix in the phase diagram [Fig.~\ref{fig:cases}(c1)].

In case 3 ($\mu=-0.92$) the spring is weak, and the upright position is a hyperbolic saddle with growth rate $(\beta|\mu|)^{1/2}=1.52\ \mathrm{s^{-1}}$. The rod departs exponentially from the vertical and rotates with a period of $6.27$~s, completing almost 16 revolutions in 100~s. Its energy exceeds the separatrix energy by only $0.2\%$, so that the orbit again follows the separatrix closely, and the landscape explored in Fig.~\ref{fig:cases}(a3) spans the full range $\pm MgL$. With $\kappa=0.08$, this case is close to the simple pendulum.

In cases 4 and 5 the larger kick, $0.4\ \mathrm{rad\,s^{-1}}$, places the energy far above the separatrix, at $4.2$ and $3.4$ times $MgL$, respectively, and the rod rotates with an angular velocity that never vanishes: the kinetic-energy parabolas of Figs.~\ref{fig:cases}(b4) and \ref{fig:cases}(b5) exclude $\dot\theta=0$, and the angle grows almost linearly in time. The landscapes explain the fine structure of these rotations. For $\mu=1$ (case 4) the potential has a local minimum at the upright position flanked by two maxima at $\theta^{*}=\pm60^{\circ}$, so that the angular velocity dips twice per revolution, whereas for $\mu=-0.4$ (case 5) the only maximum is the upright position and the velocity dips once per revolution [Figs.~\ref{fig:cases}(c4) and \ref{fig:cases}(c5)].

Figure~\ref{fig:portraits} places the five orbits on the phase cylinder, over the occupation densities of long, weakly damped Langevin trajectories at temperature $T_b=0.5\,MgL$ (with Boltzmann's constant set to one), one for each value of $\mu$, integrated with the BAOAB splitting scheme~\cite{LeimkuhlerMatthews2013}. Because the stationary density of Langevin dynamics is the Boltzmann weight $\propto e^{-H/T_b}$, the logarithmic histogram is a map of the energy landscape: it peaks in the hanging well, is depleted near the saddles, and its contours are level curves of $H$, among which the separatrix is the one through the saddles. Cases 1 and 3 run along the separatrix, case 2 is a small closed orbit in the upright well, and cases 4 and 5 are wavy rotations far above it. Taken together, the five cases display the three behaviors that a kicked upright rod can show, namely libration in the upright well, near-separatrix rotation (with critical slowing down at $\mu=0$), and fast rotation, and they confirm that the energy relative to $E_{\mathrm{sep}}$ decides among them.

\begin{figure*}
\includegraphics[width=\textwidth]{figures/fig_portraits.pdf}
\caption{Phase portraits of the five cases on the cylinder, in natural units. Panels (a)--(e) correspond to cases 1--5, with the value of $\mu$ in the panel title. Background: occupation histograms of Langevin trajectories at temperature $T_b=0.5\,MgL$, one for each value of $\mu$ (friction $0.02$, time step $0.02$, and $1.2\times10^{7}$ steps; logarithmic color scale, \textit{batlowW}~\cite{Crameri2020}). Black lines: separatrices; colored lines: orbits of the five cases, wrapped to $(-\pi,\pi]$ without spurious jumps; dots and crosses: centers and saddles.}
\label{fig:portraits}
\end{figure*}
